\section{ทฤษฎีและงานวิจัยด้านพฤติกรรมและแรงจูงใจ}

บทนี้นำเสนอการทบทวนวรรณกรรมและงานวิจัยที่เกี่ยวข้องกับการพัฒนาระบบ \textit{Pacegasus} โดยครอบคลุมทฤษฎีด้านพฤติกรรมและแรงจูงใจ หลักการออกแบบเกมมิฟิเคชันสำหรับกิจกรรมทางกาย แบบจำลองและวิธีฝึกวิ่งที่ยึดหลักฐานทางวิทยาศาสตร์ รวมถึงเทคโนโลยีและแอปพลิเคชันที่เกี่ยวข้อง เพื่อนำมาสังเคราะห์องค์ความรู้และระบุช่องว่างที่ระบบมุ่งตอบสนอง

\subsection{ทฤษฎีการกำหนดใจตนเอง}

ทฤษฎีการกำหนดใจตนเอง (Self-Determination Theory: SDT) ของ Ryan และ Deci \cite{ryan2020} อธิบายว่าแรงจูงใจที่ยั่งยืนเกิดจากการตอบสนองความต้องการพื้นฐานทางจิตใจ 3 ประการ ได้แก่ ความเป็นอิสระ (Autonomy) ความสามารถ (Competence) และความสัมพันธ์ทางสังคม (Relatedness) ความเป็นอิสระหมายถึงความรู้สึกว่าบุคคลสามารถกำหนดทิศทางและตัดสินใจเกี่ยวกับพฤติกรรมของตนเองได้ ความสามารถหมายถึงการรับรู้ว่าตนมีศักยภาพและสามารถทำสิ่งต่าง ๆ ให้สำเร็จ ส่วนความสัมพันธ์ทางสังคมหมายถึงความรู้สึกเชื่อมโยง อบอุ่น และได้รับการยอมรับจากผู้อื่น

ในการออกแบบ \textit{Pacegasus} ระบบเปิดโอกาสให้ผู้ใช้เลือกเป้าหมายการวิ่งและกำหนดตารางฝึกตามช่วงเวลาที่สะดวก เพื่อสนับสนุนความเป็นอิสระ ระบบใช้ข้อมูลสุขภาวะหลังการออกกำลังกายเพื่อปรับแผนและทำให้ผู้ใช้รับรู้ความสำเร็จอย่างต่อเนื่อง เพื่อส่งเสริมความสามารถ และมีเครือข่ายสังคมที่ผู้ใช้เพิ่มเพื่อน ชวนกันออกวิ่ง ให้กำลังใจ และจัดกิจกรรมร่วมกัน เพื่อส่งเสริมความสัมพันธ์ทางสังคม งานวิจัยยังชี้ว่าการตอบสนองความต้องการด้านความสามารถช่วยลดผลกระทบเชิงลบของความสมบูรณ์แบบที่ไม่เหมาะสม ซึ่งอาจเป็นปัจจัยที่ทำให้นักวิ่งเลิกซ้อมได้ \cite{harnois2022}

\subsection{เทคนิคการเปลี่ยนพฤติกรรม}

เทคนิคการเปลี่ยนพฤติกรรม (Behavior Change Techniques: BCTs) คือกลยุทธ์เชิงโครงสร้างที่มีหลักฐานเชิงประจักษ์สนับสนุนว่าสามารถปรับเปลี่ยนพฤติกรรมได้อย่างมีประสิทธิผล การทบทวนวรรณกรรมของ Milne-Ives และคณะ \cite{milneives2023} พบว่า BCTs ที่ช่วยเพิ่มการมีส่วนร่วมในแอปพลิเคชันสุขภาพ ได้แก่ การตั้งเป้าหมาย (Goal Setting) การให้ข้อมูลป้อนกลับและติดตามผล (Feedback and Monitoring) และการเปรียบเทียบทางสังคม (Social Comparison)

ระบบ \textit{Pacegasus} ให้ผู้ใช้กำหนดเป้าหมายหลัก เช่น การวิ่งมินิมาราธอนภายใน 3 เดือน แล้วแปลงเป็นแผนฝึกรายสัปดาห์และรายวันโดยอัตโนมัติ หลังวิ่ง ระบบแสดงเพซ ระยะทาง และข้อมูลความก้าวหน้าเปรียบเทียบกับสัปดาห์ก่อน พร้อมช่องบันทึกความรู้สึกเพื่อนำไปปรับแผน ส่วนการเปรียบเทียบทางสังคมใช้กระดานผู้นำในกลุ่มเพื่อน เพื่อสร้างการแข่งขันเชิงบวกโดยไม่ก่อให้เกิดแรงกดดันเกินควร

\subsection{แรงจูงใจและการสนับสนุนทางสังคม}

พฤติกรรมการออกกำลังกายได้รับอิทธิพลจากบุคคลรอบข้างอย่างมีนัยสำคัญ Golaszewski และคณะ \cite{golaszewski2022} พบว่าการออกกำลังกายแบบกลุ่มส่งเสริมการสนับสนุนจากเพื่อนร่วมกิจกรรม การสร้างอัตลักษณ์ของการเป็นนักออกกำลังกาย และปริมาณการออกกำลังกายในระดับปานกลางถึงหนัก เมื่อผู้ใช้มองว่าการวิ่งเป็นส่วนหนึ่งของตัวตน พฤติกรรมดังกล่าวย่อมมีแนวโน้มยั่งยืนมากขึ้น

ด้วยเหตุนี้ ระบบ \textit{Pacegasus} จึงมีระบบกลุ่ม (Club) ที่สนับสนุนการสร้างกลุ่ม การกำหนดบทบาทสมาชิก โปรไฟล์ทีม โลโก้ สีประจำทีม และสโลแกน เพื่อสร้างความรู้สึกเป็นเจ้าของและอัตลักษณ์ร่วมของผู้ใช้

\section{หลักการออกแบบเกมมิฟิเคชันสำหรับกิจกรรมทางกาย}

เกมมิฟิเคชันหมายถึงการนำองค์ประกอบและกลไกของเกมมาประยุกต์ใช้ในบริบทที่ไม่ใช่เกม \cite{fu2024} ระบบ \textit{Pacegasus} นำหลักการดังกล่าวมาประยุกต์ใช้ ดังนี้

\subsection{ระบบเป้าหมายและความก้าวหน้า}

ระบบกำหนดเป้าหมายทั้งระยะสั้นและระยะยาว พร้อมแสดงความก้าวหน้าแบบต่อเนื่อง เป้าหมายหลัก เช่น การเตรียมตัวสำหรับการแข่งขัน ถูกย่อยเป็นเป้าหมายรายสัปดาห์และรายวัน เพื่อให้ผู้ใช้เห็นพัฒนาการอย่างชัดเจนและเกิดความรู้สึกมีความสามารถ

\subsection{ระบบการให้รางวัลและสิ่งจูงใจ}

ระบบรางวัลประกอบด้วยคะแนนสะสมจากกิจกรรมรายวัน เหรียญตราสำหรับภารกิจตามเงื่อนไข และระดับเลเวลที่ปลดล็อกฟีเจอร์ใหม่เมื่อผู้ใช้สะสมคะแนนถึงเกณฑ์ที่กำหนด การให้รางวัลได้รับการออกแบบให้สอดคล้องกับระดับความก้าวหน้าของผู้ใช้ เพื่อคงแรงจูงใจในการใช้งาน \cite{milneives2023}

\subsection{ความท้าทายที่ปรับตามระดับผู้ใช้}

ภารกิจและแผนการฝึกในระบบปรับระดับความยากตามสมรรถภาพของผู้ใช้แต่ละคนผ่าน Adaptive Training Engine เพื่อคงความสมดุลระหว่างความท้าทายและทักษะของผู้ใช้ ซึ่งสนับสนุนให้เกิดความพึงพอใจและแรงจูงใจในการออกกำลังกายอย่างต่อเนื่อง

\subsection{องค์ประกอบทางสังคมและเนื้อเรื่อง}

ระบบรองรับการสร้างและเข้าร่วมกลุ่ม การแข่งขันบนกระดานผู้นำในกลุ่มเพื่อน และการให้กำลังใจซึ่งกันและกัน ซึ่งเป็นคุณลักษณะที่ช่วยเพิ่มระดับกิจกรรมทางกาย \cite{tong2018} นอกจากนี้ Story Mode เชื่อมโยงกิจกรรมการวิ่งเข้ากับเนื้อเรื่องต่อเนื่อง โดยการวิ่งแต่ละครั้งสามารถปลดล็อกตอนใหม่ ทำให้ผู้ใช้มีเหตุผลในการกลับมาใช้งานอย่างสม่ำเสมอ พร้อมข้อมูลผลลัพธ์ทันทีหลังจบกิจกรรม

\section{แบบจำลองและวิธีฝึกวิ่งที่ยึดหลักฐานทางวิทยาศาสตร์}

\subsection{การใช้ Machine Learning ในการวางแผนการฝึกซ้อม}

Kerimov \cite{kerimov2025} แสดงให้เห็นว่าแบบจำลอง Machine Learning สามารถทำนายสมรรถภาพและความพร้อมในการฝึกของนักกีฬาได้จากข้อมูลที่ผู้ใช้รายงานเอง เช่น ความเหนื่อยล้า คุณภาพการนอนหลับ ความเจ็บปวดของกล้ามเนื้อ ความเครียด และอารมณ์โดยรวม แนวคิดนี้เป็นพื้นฐานของฟีเจอร์ Daily Wellness Check-in ในระบบ \textit{Pacegasus} ซึ่งให้ผู้ใช้ประเมินสุขภาวะของตนเองในเวลาสั้น ๆ แล้วใช้ข้อมูลเป็นส่วนหนึ่งของการปรับโปรแกรมฝึกประจำวัน

\subsection{Session-RPE}

Session Rating of Perceived Exertion (Session-RPE หรือ sRPE) คำนวณจากระยะเวลาที่ออกกำลังกาย (นาที) คูณด้วยระดับความเหนื่อยล้าที่รับรู้ในสเกล 1--10 งานวิจัยของ Dai, Yan และ Bi \cite{dai2025} ยืนยันว่าสามารถใช้เป็นมาตรวัดภาระการฝึกที่น่าเชื่อถือ เข้าถึงง่าย และไม่ต้องอาศัยอุปกรณ์เฉพาะทาง ระบบ \textit{Pacegasus} จึงใช้ Post-Run Log ให้ผู้ใช้ประเมินระดับความเหนื่อยหลังวิ่ง เพื่อแสดงคะแนนภาระการฝึกของวันนั้น

\subsection{Acute-to-Chronic Workload Ratio}

Acute-to-Chronic Workload Ratio (ACWR) เป็นเครื่องมือสำหรับบริหารจัดการความเสี่ยงต่อการบาดเจ็บ การทบทวนวรรณกรรมและการวิเคราะห์อภิมานของ Qin, Li และ Chen \cite{qin2025} พบว่าช่วงค่าที่เหมาะสมอยู่ระหว่าง 0.8--1.3 และเมื่อค่าสูงกว่า 1.5 ความเสี่ยงต่อการบาดเจ็บจะเพิ่มขึ้น ระบบ \textit{Pacegasus} คำนวณ ACWR จากข้อมูล sRPE ที่สะสมไว้ แล้วแสดงผลเป็นดัชนีความเสี่ยง 3 ระดับ ได้แก่ สีเขียว สีเหลือง และสีแดง พร้อมแจ้งเตือนเชิงป้องกันเมื่อค่าเพิ่มขึ้นอย่างรวดเร็ว

\section{เทคโนโลยีและแอปพลิเคชันที่เกี่ยวข้อง}

เพื่อกำหนดตำแหน่งของระบบ \textit{Pacegasus} ในระบบนิเวศของแอปพลิเคชันส่งเสริมสุขภาพ ตารางที่ \ref{tab:running-app-comparison} เปรียบเทียบคุณลักษณะสำคัญของแอปพลิเคชันที่เกี่ยวข้อง

\begin{table}[htbp]
\centering
\caption{การเปรียบเทียบคุณลักษณะของแอปพลิเคชันวิ่งที่เกี่ยวข้อง}
\label{tab:running-app-comparison}
\small
\begin{tabular}{p{0.30\textwidth}cccc}
\hline
\textbf{คุณลักษณะ} & \textbf{Strava} & \textbf{Nike Run Club} & \textbf{Garmin Connect} & \textbf{Pacegasus} \\
\hline
การติดตาม GPS & มี & มี & มี & มี \\
Adaptive Training Engine & ไม่มี & บางส่วน & บางส่วน & เต็มรูปแบบ \\
Gamification & จำกัด & มี & จำกัด & เต็มรูปแบบ \\
Story Mode & ไม่มี & ไม่มี & ไม่มี & มี \\
Social Running Community & มี & จำกัด & ไม่มี & มี \\
ต้องการอุปกรณ์เสริม & ไม่ต้องการ & ไม่ต้องการ & Garmin Watch & ไม่ต้องการ \\
Session-RPE / ACWR & ไม่มี & ไม่มี & บางส่วน & มี \\
Daily Wellness Check-in & ไม่มี & ไม่มี & มี & มี \\
\hline
\end{tabular}
\end{table}

จากการวิเคราะห์พบว่าแอปพลิเคชันในตลาดยังขาดการบูรณาการ Adaptive Training Engine เกมมิฟิเคชันเชิงเนื้อเรื่อง และระบบสังคมที่ออกแบบบนพื้นฐานของ SDT อย่างครบถ้วน อีกทั้งแอปพลิเคชันที่มีฟีเจอร์ครบถ้วนมักต้องพึ่งพาอุปกรณ์สวมใส่ราคาแพง ซึ่งเป็นอุปสรรคต่อผู้ใช้ทั่วไป ช่องว่างดังกล่าวคือสิ่งที่ \textit{Pacegasus} มุ่งตอบสนอง

\section{สรุปกรอบแนวคิดการวิจัย}

จากการทบทวนวรรณกรรม SDT เป็นกรอบหลักสำหรับอธิบายแรงจูงใจภายใน BCTs เป็นชุดกลยุทธ์สำหรับการเปลี่ยนพฤติกรรม เกมมิฟิเคชันเป็นกลไกเพิ่มการมีส่วนร่วมและการกลับมาใช้งาน หลักการวิทยาศาสตร์การกีฬา ได้แก่ Session-RPE, ACWR และ Machine Learning สนับสนุนความปลอดภัยและประสิทธิภาพของแผนการฝึก ขณะที่คุณลักษณะทางสังคมช่วยสร้างแรงจูงใจและอัตลักษณ์ของการเป็นนักวิ่ง การบูรณาการทั้งห้าด้านเป็นจุดแข็งสำคัญของระบบ \textit{Pacegasus} เมื่อเปรียบเทียบกับแอปพลิเคชันอื่นในตลาด
